\documentclass[10pt,a4paper]{amsart}
\usepackage[margin=19mm]{geometry}
\usepackage{amsmath,amssymb,mathtools}
\usepackage[scr=rsfs,cal=euler]{mathalfa}
\usepackage{xfrac}
\usepackage[unicode,hidelinks,bookmarks=false]{hyperref}
\newcommand{\ie}{i.e.\ }
\newcommand{\cf}{cf.\ }
\newcommand{\resp}{respectively\ }
\newcommand{\Subsection}{\subsection}
\providecommand{\nobreakdash}{\nobreak\hbox{-}}
\setlength{\emergencystretch}{2em}
\multlinegap0pt
\usepackage{amssymb}
\usepackage{float}
\newtheorem{lemma}{Lemma}
\DeclareMathOperator{\End}{End}
\DeclareMathOperator{\Hom}{Hom}
\DeclareMathOperator{\Soc}{Soc}
\DeclareMathOperator{\rad}{rad}
\title{On the minimal dimension of maximal commutative subalgebras of $M_6(k)$}
\author{Małgorzata Nowak-Kępczyk}
\date{Source: arXiv:2604.23322v1 (CC BY 4.0)}
\begin{document}
\maketitle

\noindent\emph{Source and licence.} On the minimal dimension of maximal commutative subalgebras of $M_6(k)$. Version arXiv:2604.23322v1 (CC BY 4.0), distributed by arXiv under the Creative Commons Attribution 4.0 International licence, http://creativecommons.org/licenses/by/4.0/, as recorded in the arXiv licence metadata for this exact version and shown on the abstract page https://arxiv.org/abs/2604.23322v1. Full licence notice: \texttt{licenses/arxiv-2604.23322-CC-BY-4.0.txt}.\par\smallskip

\noindent\textit{Editorial scope.} Counted unit: the article's lemma lem:class16 (source0.tex lines 474--490). The article's complete original proof of it is delivered with it (source0.tex lines 492--678, one \texttt{proof} environment of 187 lines). No mathematics is written, repaired or re-derived: the statement and the proof are the article's own lines, in the article's own notation, and no retained line is substituted or re-flowed.  No further source-local context is needed: every cross-reference of every retained line resolves inside the two delivered spans.  Nothing was omitted from any retained span, so the package carries no omission record.  No retained line contains a citation, so the wrapper carries no bibliography and no bibliography entry.  No retained line contains a figure environment, an image-inclusion command or drawing code, so no image is delivered and the package declares no asset dependency.  Editorial formatting only, in addition: an AMS \texttt{amsart} preamble that restates the article's own declarations and macros used by the retained lines, namely the environments \texttt{lemma} on separate counters, together with the standard packages those lines need.  The retained environments print in this document as Lemma 1 in order of appearance, whereas the article prints the same items with the numbering of its own full document.  The complete native source of the article as distributed in the arXiv e-print for this exact version is bundled unchanged as \texttt{sources/199155/source0.tex}.
\begin{lemma}\label{lem:class16}
	Let
	\[
	A \cong k[x,y,z]/(x^3,\ y^2,\ z^2,\ xy,\ xz,\ yz)
	\]
	be the local algebra of class {\rm 16}, and let $V=k^6$ be a faithful $A$-module such that
	\[
	\dim(V/JV,\ JV/J^2V,\ J^2V)=(2,3,1),
	\]
	where $J=\rad(A)=(x,y,z)$ and $J^2=(x^2)$.
	
	Then
	\[
	\dim \End_A(V)\ge 6.
	\]
	In particular, the image of $A$ in $M_6(k)$ is not a maximal commutative subalgebra.
\end{lemma}

\begin{proof}
	Write
	\[
	w:=x^2.
	\]
	Then
	\[
	J=\langle x,y,z,w\rangle,\qquad J^2=\langle w\rangle,
	\]
	and the defining relations imply
	\[
	yJ=0,\qquad zJ=0,\qquad wJ=0.
	\]
	
	Let
	\[
	V=V_0\oplus V_1\oplus V_2,
	\qquad \dim V_0=2,\ \dim V_1=3,\ \dim V_2=1,
	\]
	be a decomposition compatible with the filtration
	\[
	V\supset JV\supset J^2V\supset 0.
	\]
	Each $u\in J$ has block form
	\[
	u=
	\begin{pmatrix}
		0&0&0\\
		L_u&0&0\\
		M_u&N_u&0
	\end{pmatrix},
	\]
	where
	\[
	L_u:V_0\to V_1,\qquad N_u:V_1\to V_2,\qquad M_u:V_0\to V_2.
	\]
	
	Since $JV/J^2V=V_1$, the images of $L_x,L_y,L_z$ span $V_1$. Because $yJ=0$, we have
	\[
	N_yL_x=N_yL_y=N_yL_z=0,
	\]
	hence $N_y=0$. Similarly, $N_z=0$.
	
	On the other hand, $w=x^2$ acts nontrivially on $V$, since $V$ is faithful and $w\neq 0$ in $A$.
	Since the operator $x^2$ is represented by the block
	\[
	N_xL_x:V_0\to V_2,
	\]
	we have
	\[
	N_xL_x\neq 0.
	\]
	In particular, $N_x\neq 0$.
	
	Now consider the subspace
	\[
	U:=\operatorname{Im}L_y+\operatorname{Im}L_z \subset V_1.
	\]
	Because $xy=xz=0$, we have
	\[
	N_xL_y=N_xL_z=0,
	\]
	so
	\[
	U\subset \ker N_x.
	\]
	Since $N_x\neq 0$ and $\dim V_1=3$, it follows that $\dim\ker N_x=2$.
	
	We distinguish two cases.
	
	\medskip
	\noindent
	\textbf{Case 1:} $\dim U=2$.
	
	Then $U=\ker N_x$. Since $yJ=0$ and $zJ=0$, the images of the operators $y$ and $z$
	are contained in the socle
	\[
	\Soc_A(V):=\{v\in V:\ Jv=0\}.
	\]
	Moreover, $U\subset \Soc_A(V)$ and $V_2=J^2V\subset \Soc_A(V)$, so
	\[
	\dim \Soc_A(V)\ge 3.
	\]
	Therefore every map
	\[
	\varphi:V/JV\to \Soc_A(V)
	\]
	induces an $A$-endomorphism of $V$, and we obtain an embedding
	\[
	\Hom_k(V/JV,\Soc_A(V))\hookrightarrow \End_A(V).
	\]
	Since
	\[
	\dim(V/JV)=2,\qquad \dim\Soc_A(V)\ge 3,
	\]
	it follows that
	\[
	\dim\End_A(V)\ge 2\cdot 3=6.
	\]
	Adding scalar endomorphisms, we obtain in fact
	\[
	\dim\End_A(V)\ge 7.
	\]
	
	\medskip
	\noindent
	\textbf{Case 2:} $\dim U=1$.
	
	Since $L_y$ and $L_z$ are linearly independent modulo $J^2$, they are linearly independent as maps
	$V_0\to V_1$. Thus, after changing bases in $V_0$ and $V_1$, we may assume that
	\[
	L_y=
	\begin{pmatrix}
		1&0\\
		0&0\\
		0&0
	\end{pmatrix},
	\qquad
	L_z=
	\begin{pmatrix}
		0&1\\
		0&0\\
		0&0
	\end{pmatrix}.
	\]
	Since $N_x\neq 0$ and $N_xL_y=N_xL_z=0$, we may also assume that
	\[
	N_x=(0\ 0\ 1).
	\]
	Finally, since $N_xL_x\neq 0$, after changing the basis in a complement of
	$\operatorname{Im}L_y+\operatorname{Im}L_z$, we may write
	\[
	L_x=
	\begin{pmatrix}
		0&0\\
		1&0\\
		0&1
	\end{pmatrix}.
	\]
	
	Let
	\[
	X=
	\begin{pmatrix}
		P&0&0\\
		U'&Q&0\\
		W&T&r
	\end{pmatrix}\in \End_A(V).
	\]
	The relations $Xu=uX$ for $u=x,y,z$ yield
	\[
	rN_x=N_xQ,\qquad QL_u=L_uP\quad (u=x,y,z),
	\]
	together with
	\[
	TL_u+rM_u=N_uU'+M_uP\quad (u=x,y,z).
	\]
	
	A straightforward computation with the above normal forms shows that these equations force
	\[
	P=rI_2,\qquad Q=rI_3,
	\]
	and impose only the additional conditions
	\[
	t_1=0,\qquad u_{31}=t_2,\qquad u_{32}=t_3,
	\]
	while the remaining entries of $U'$, all entries of $W\in\Hom(V_0,V_2)$, and the parameters
	$t_2,t_3,r$ are free.
	
	Thus $\End_A(V)$ contains at least the following free parameters:
	\[
	r \quad (1),\qquad W\in \Hom(V_0,V_2)\quad (2),
	\]
	\[
	\text{the first two rows of }U' \quad (4),\qquad t_2,t_3 \quad (2).
	\]
	Hence
	\[
	\dim\End_A(V)\ge 1+2+4+2=9.
	\]
	
	In both cases,
	\[
	\dim\End_A(V)>\dim A=5.
	\]
	Therefore the image of $A$ in $M_6(k)$ cannot be maximal commutative.
\end{proof}

\end{document}
